
\section{诺维科夫覆盖与作用量泛函}
\label{sec:2-2}

设 $(L_0,L_1)$ 是 $(M,\omega)$ 中一对相互横截的紧拉格朗日子流形。这里要指出，目前我们并不假设这些拉格朗日子流形是连通的。

考虑道路空间
\[
  \Omega=\Omega(L_0,L_1)
  =\{\ell:[0,1]\longrightarrow M
    \mid \ell(0)\in L_0,\ \ell(1)\in L_1\}.
\]
在这个空间上，定义作用 $1$-形式 $\alpha$：
\begin{equation}
  \alpha(\ell)(\xi)
  =\int_0^1\omega\bigl(\dot\ell(t),\xi(t)\bigr)\,dt,
  \tag{2.2.1}\label{eq:2-2-1}
\end{equation}
其中每个 $\xi\in T_\ell\Omega$。利用 $(L_0,L_1)$ 的拉格朗日性质，直接计算可知这个形式是闭的（见命题 2.2.8）。弗洛尔同调理论是闭 $1$-形式的诺维科夫--莫尔斯理论的一种半无限维版本。

注意，$\Omega(L_0,L_1)$ 并不连通，而是有可数多个连通分支。研究 $\Omega(L_0,L_1)$ 的某个固定连通分支时，我们选取一条基准道路，记为 $\ell_0$，以指定这个分支。相应分支记作
$\Omega(L_0,L_1;\ell_0)\subset\Omega(L_0,L_1)$。基准道路 $\ell_0$ 的始点和终点
\[
  x_0=\ell_0(0)\in L_0,
  \qquad
  x_1=\ell_0(1)\in L_1
\]
自动分别从 $L_0$ 与 $L_1$ 中选出一个连通分支。于是，$\Omega(L_0,L_1;\ell_0)$ 是相应两个连通分支之间的道路空间的一个子空间。因此，除非另有说明，从现在起总假设 $L_0,L_1$ 连通。

接下来描述一个覆盖空间，称为 $\Omega(L_0,L_1)$ 的分支 $\Omega(L_0,L_1;\ell_0)$ 的诺维科夫覆盖\liftIndex{I063}。先从 $\Omega(L_0,L_1;\ell_0)$ 的万有覆盖空间开始。考虑所有对子 $(\ell,w)$，其中 $\ell\in\Omega(L_0,L_1)$，且映射 $w:[0,1]^2\to M$ 满足边界条件
\begin{equation}
  \begin{cases}
    w(0,\mathord\cdot)=\ell_0,
      &w(1,\mathord\cdot)=\ell,\\
    w(s,0)\in L_0,
      &w(s,1)\in L_1\quad\text{对所有 }s\in[0,1].
  \end{cases}
  \tag{2.2.2}\label{eq:2-2-2}
\end{equation}
把 $w$ 看成 $\Omega(L_0,L_1;\ell_0)$ 中的连续道路 $s\mapsto w(s,\mathord\cdot)$，其两端为 $\ell_0$ 与 $\ell$；则 $\Omega(L_0,L_1;\ell_0)$ 的万有覆盖空间在 $\ell$ 上的纤维，可以由 $w$ 相对于两端 $s=0,1$ 的道路同伦类集合表示。

\subsection{\texorpdfstring{$\Gamma$}{Gamma}-等价}
\label{sec:2-2-1}

为了实现诺维科夫的莫尔斯理论思想\textbf{[Nov81]}，关键是要有 $\Omega(L_0,L_1;\ell_0)$ 的某个阿贝尔覆盖。就此而言，万有覆盖空间过于庞大。现在用另一个比同伦更弱的等价关系，对 $\Omega(L_0,L_1;\ell_0)$ 中的道路空间取商，从而定义 $\Omega(L_0,L_1;\ell_0)$ 的一个较小覆盖空间。它类似于霍费尔--萨拉蒙\textbf{[HoSa95]} 和第四署名作者小野薫\textbf{[Ono95]} 所引入的可缩环路空间的诺维科夫覆盖空间。由构造可见，这个覆盖空间的覆盖变换群是阿贝尔群。

给定来自 $\Omega(L_0,L_1;\ell_0)$ 的两个对子 $(\ell,w)$ 与 $(\ell,w')$ 时，拼接
\[
  \overline w\#w':[0,1]\times[0,1]\longrightarrow M
\]
定义了一个环路 $c:S^1\to\Omega(L_0,L_1;\ell_0)$。可以把这个环路看成一个映射
\[
  C:S^1\times[0,1]\longrightarrow M,
\]
它满足边界条件
$C(s,0)\in L_0$、$C(s,1)\in L_1$。显然，$C$ 的辛面积记为
\begin{equation}
  I_\omega(c)=\int_C\omega.
  \tag{2.2.3}\label{eq:2-2-3}
\end{equation}
它只依赖于满足上述 $L_0/L_1$ 边界条件的 $C$ 的同伦类，因而在
$\pi_1(\Omega(L_0,L_1;\ell_0))$ 上定义了一个同态，仍记为
\[
  I_\omega:\pi_1\bigl(\Omega(L_0,L_1;\ell_0)\bigr)
  \longrightarrow\mathbb R.
\]

其次，对满足上述 $L_0/L_1$ 边界条件的映射
$C:S^1\times[0,1]\to M$，配置辛丛对 $(\mathcal V_C,\lambda_C)$：
\[
  \mathcal V_C=C^*TM,
  \qquad
  \lambda_C=c_0^*TL_0\amalg c_1^*TL_1,
\]
其中 $c_i:S^1\to L_i$ 由
$c_i(s)=C(s,i)$（$i=0,1$）给出。由此可以定义另一个同态
\[
  I_\mu:\pi_1\bigl(\Omega(L_0,L_1;\ell_0)\bigr)
  \longrightarrow\mathbb Z,
  \qquad
  I_\mu(c)=\mu(\mathcal V_C,\lambda_C),
\]
其中 $\mu(\mathcal V_C,\lambda_C)$ 是丛对 $(\mathcal V_C,\lambda_C)$ 的马斯洛夫指标\liftIndex{I054}。

利用同态 $I_\mu$ 和 $I_\omega$，在所有满足式~\eqref{eq:2-2-2} 的对子 $(\ell,w)$ 之集合上定义等价关系 $\sim$。对两个这样的填充 $w,w'$，以 $\overline w\#w'$ 表示二者沿 $\ell$ 的拼接；它在 $\Omega(L_0,L_1;\ell_0)$ 中定义了一条以 $\ell_0$ 为基点的环路。我们借用塞德尔在哈密顿微分同胚的弗洛尔同调语境中所用的术语\textbf{[Sei97]}，引入 $\Gamma$-等价这一概念。

\begin{definition}{2.2.4}\label{def:2-2-4}
若
\[
  I_\omega(\overline w\#w')=0
  =I_\mu(\overline w\#w'),
\]
则称 $(\ell,w)$ 与 $(\ell,w')$ 是 $\Gamma$-等价的，并记作
$(\ell,w)\sim(\ell,w')$。以
\[
  \widetilde\Omega(L_0,L_1;\ell_0)
\]
表示等价类 $[\ell,w]$ 的集合，并称之为诺维科夫覆盖空间。
\end{definition}

$\ell_0\in\Omega(L_0,L_1;\ell_0)$ 有一个到
$\widetilde\Omega(L_0,L_1;\ell_0)$ 的典范提升：它就是
\[
  [\ell_0,\widetilde\ell_0]
  \in\widetilde\Omega(L_0,L_1;\ell_0),
\]
其中 $\widetilde\ell_0:[0,1]^2\to M$ 是映射
\[
  \widetilde\ell_0(s,t)=\ell_0(t).
\]
这样，$\widetilde\Omega(L_0,L_1;\ell_0)$ 也有一个自然基点；我们在记号中省略它。

以 $\Pi(L_0,L_1;\ell_0)$ 表示覆盖
\[
  \widetilde\Omega(L_0,L_1;\ell_0)
  \longrightarrow\Omega(L_0,L_1;\ell_0)
\]
的覆盖变换群。不难看出，$\Pi(L_0,L_1;\ell_0)$ 的同构类只依赖于包含 $\ell_0$ 的连通分支。

两个同态 $I_\omega$ 与 $I_\mu$ 下降为同态
\[
  E:\Pi(L_0,L_1;\ell_0)\longrightarrow\mathbb R,
  \qquad
  \mu:\Pi(L_0,L_1;\ell_0)\longrightarrow\mathbb Z,
\]
其定义为
\[
  E(g)=I_\omega[C],
  \qquad
  \mu(g)=I_\mu[C],
\]
其中 $C:S^1\times[0,1]\to M$ 是代表类
$g\in\Pi(L_0,L_1;\ell_0)$ 的任一映射。

\begin{corollary}{2.2.5}\label{cor:2-2-5}
群 $\Pi(L_0,L_1;\ell_0)$ 是阿贝尔群。
\end{corollary}

\begin{proof}
由 $\Pi(L_0,L_1;\ell_0)$ 的定义，映射
\[
  E\times\mu:\Pi(L_0,L_1;\ell_0)
  \longrightarrow\mathbb R\times\mathbb Z
\]
是单射群同态。由于 $\mathbb R\times\mathbb Z$ 是阿贝尔群，故
$\Pi(L_0,L_1;\ell_0)$ 也是阿贝尔群。
\end{proof}

现在把群环 $R[\Pi(L_0,L_1;\ell_0)]$ 完备化，定义与阿贝尔覆盖
\[
  \widetilde\Omega(L_0,L_1;\ell_0)
  \longrightarrow\Omega(L_0,L_1;\ell_0)
\]
相伴的诺维科夫环\liftIndex{I064} $\Lambda(L_0,L_1;\ell_0)$。这里 $R$ 是有幺交换环。

\begin{definition}{2.2.6}\label{def:2-2-6}
$\Lambda_k^R(L_0,L_1;\ell_0)$ 表示所有无穷和
\[
  \sum_{\substack{g\in\Pi(L_0,L_1;\ell_0)\\\mu(g)=k}}
  a_g[g]
\]
所成的集合，其中 $a_g\in R$，并且满足下列条件：
\begin{equation}
  \text{对每个 }C\in\mathbb R,\quad
  \{g\in\Pi(L_0,L_1;\ell_0)
    \mid E(g)\leq C,\ a_g\neq0\}
  \text{ 是有限集。}
  \tag{2.2.7}\label{eq:2-2-7}
\end{equation}
令
\[
  \Lambda^R(L_0,L_1;\ell_0)
  =\bigoplus_k\Lambda_k^R(L_0,L_1;\ell_0).
\]
\end{definition}

我们主要使用 $R=\mathbb Q$。当 $R=\mathbb Q$ 时，以
$\Lambda(L_0,L_1;\ell_0)$ 代替
$\Lambda^{\mathbb Q}(L_0,L_1;\ell_0)$。

$\Lambda^R(L_0,L_1;\ell_0)$ 上的环结构由卷积积定义：
\[
  \left(\sum_{g\in\Pi(L_0,L_1;\ell_0)}a_g[g]\right)
  \left(\sum_{g\in\Pi(L_0,L_1;\ell_0)}b_g[g]\right)
  =\sum_{g_1,g_2\in\Pi(L_0,L_1;\ell_0)}
    a_{g_1}b_{g_2}[g_1g_2].
\]
不难看出，右端确实是 $\Lambda^R(L_0,L_1;\ell_0)$ 的一个元素，即满足其定义中的有限性条件~\eqref{eq:2-2-7}。因此，在这个乘法下，
\[
  \Lambda^R(L_0,L_1;\ell_0)
  =\bigoplus_k\Lambda_k^R(L_0,L_1;\ell_0)
\]
成为一个分次环。我们称它为与对子 $(L_0,L_1)$ 及包含 $\ell_0$ 的连通分支相伴的诺维科夫环。

\subsection{作用量泛函与马斯洛夫--莫尔斯指标}
\label{sec:2-2-2}
\liftIndex{I056}

\begin{samepage}
对给定的对子 $(\ell,w)$，定义作用量泛函\liftIndex{I039}
\[
  \mathcal A:\widetilde\Omega(L_0,L_1;\ell_0)
  \longrightarrow\mathbb R
\]
为
\[
  \mathcal A(\ell,w)=\int w^*\omega.
\]
\end{samepage}
由 $\Pi(L_0,L_1;\ell_0)$ 的定义可知，这个积分只依赖于 $\Gamma$-等价类 $[\ell,w]$，因而下降为覆盖空间 $\widetilde\Omega(L_0,L_1;\ell_0)$ 上一个良定义的泛函。直接计算可得下面的命题；证明从略。

\begin{proposition}{2.2.8}\label{prop:2-2-8}
设
\[
  \pi:\widetilde\Omega(L_0,L_1;\ell_0)
  \longrightarrow\Omega(L_0,L_1;\ell_0)
\]
为该 $\Gamma$-覆盖的投影，$\alpha$ 为 $\Omega(L_0,L_1;\ell_0)$ 上的作用 $1$-形式。则
\[
  d\mathcal A=-\pi^*\alpha.
\]
\end{proposition}

由这个命题和 $\alpha$ 的定义立即得到 $\mathcal A$ 的临界点集合的如下刻画。

\begin{corollary}{2.2.9}\label{cor:2-2-9}
$\mathcal A$ 的临界点集合 $\operatorname{Cr}(L_0,L_1;\ell_0)$ 由形如 $[\ell_p,w]$ 的对子组成，其中 $\ell_p$ 是取值为某个 $p\in L_0\cap L_1$ 的常值道路，$w$ 如式~\eqref{eq:2-2-2} 所示。集合 $\operatorname{Cr}(L_0,L_1;\ell_0)$ 在 $\Pi(L_0,L_1;\ell_0)$ 的作用下不变，因而构成 $L_0\cap L_1$ 的某个子集上的主丛，其纤维同构于 $\Pi(L_0,L_1;\ell_0)$。
\end{corollary}

令
\[
  \operatorname{Cr}(L_0,L_1)
  =\bigcup_{\ell_{0,i}}
    \operatorname{Cr}(L_0,L_1;\ell_{0,i}),
\]
其中 $\ell_{0,i}$ 遍历 $\Omega(L_0,L_1)$ 各连通分支的基点集合。

接下来给 $\mathcal A$ 的每个临界点赋予一个绝对莫尔斯指标。一般而言，这并非平凡之事，因为 $\mathcal A$ 在任一临界点处显然的莫尔斯指标都是无穷。为此，将使用与临界点
$[\ell_p,w]\in\operatorname{Cr}(L_0,L_1;\ell_0)$ 自然相伴的某个丛对的马斯洛夫指标。

我们把 $[\ell_p,w]$ 的这个莫尔斯指标称为该临界点的马斯洛夫--莫尔斯指标（相对于基准道路 $\ell_0$）。$\mu([\ell_p,w])$ 的定义与 $\mathcal A$ 的定义有几分相似。不过，为定义它，还需要固定 $\ell_0^*\Lambda(M)$ 的一个截面 $\lambda^0$，使得
\[
  \lambda^0(0)=T_{\ell_0(0)}L_0,
  \qquad
  \lambda^0(1)=T_{\ell_0(1)}L_1.
\]
这里，$\Lambda(M)$ 是 $TM$ 的拉格朗日格拉斯曼流形丛\liftIndex{I055}：
\[
  \Lambda(M)=\bigcup_{p\in M}\Lambda(T_pM),
\]
其中 $\Lambda(T_pM)$ 是辛向量空间 $(T_pM,\omega_p)$ 的拉格朗日子空间集合。令
\[
  \Lambda^{\mathrm{ori}}(M)
  =\bigcup_{p\in M}\Lambda^{\mathrm{ori}}(T_pM)
\]
为定向拉格朗日格拉斯曼流形丛；它是 $\Lambda(M)$ 的二重覆盖。当 $L_0,L_1$ 有定向时，假设 $\lambda^0$ 是 $\Lambda^{\mathrm{ori}}(M)$ 的截面，并在两个端点处保持定向。

取元素
\[
  [\ell_p,w]\in\operatorname{Cr}(L_0,L_1;\ell_0)
  \subset\widetilde\Omega(L_0,L_1;\ell_0),
\]
其投影对应于交点 $p\in L_0\cap L_1$。

利用命题 2.1.3，在方形 $[0,1]^2$ 上为它配置一个辛丛对 $(\mathcal V_w,\lambda_w)$；这个丛对在同伦意义下唯一确定。先取
\[
  \mathcal V_w=w^*TM.
\]
为定义 $\lambda_w$，选取道路
$\alpha^p:[0,1]\to\Lambda(T_pM,\omega_p)$，使得
\begin{equation}
  \begin{cases}
    \alpha^p(0)=T_pL_0,
      \quad \alpha^p(1)=T_pL_1\subset T_pM,\\
    \alpha^p(t)\pitchfork T_pL_0
      \quad\text{对所有 }0<t\leq1,\\
    \alpha^p(t)\in U_0(T_pL_0)
      \quad\text{当 $t$ 很小时},
  \end{cases}
  \tag{2.2.10}\label{eq:2-2-10}
\end{equation}
其中 $U_0(T_pL_0)$ 如式 (2.1.2) 所定义。命题 2.1.3 蕴含，$\alpha^p$ 的任意两个这种选取彼此同伦。

再由下式定义 $\mathcal V_w|_{\partial[0,1]^2}$ 的连续拉格朗日子丛
$\lambda_w\to\partial[0,1]^2$；它在 $\partial[0,1]^2$ 每一点处的纤维为
\translatornote{原书在 (2.2.10) 漏去横截条件的范围 $0<t\leq1$，又在 (2.2.11) 把 $\mathcal V_w$ 漏写下标，并将一参数截面误写成 $\lambda^0(0,t)$；译文依命题 2.1.3 及本段定义补正。}
\begin{equation}
  \begin{cases}
    \lambda_w(s,0)=T_{w(s,0)}L_0,
      &\lambda_w(1,t)=\alpha^p(t),\\
    \lambda_w(s,1)=T_{w(s,1)}L_1,
      &\lambda_w(0,t)=\lambda^0(t).
  \end{cases}
  \tag{2.2.11}\label{eq:2-2-11}
\end{equation}
命题 2.1.3 还蕴含，上述构造的丛对之同伦型也与 $\alpha^p$ 的选取无关。

\begin{definition}{2.2.12}\label{def:2-2-12}
定义 $[\ell_p,w]$ 的马斯洛夫--莫尔斯指标（相对于 $\lambda^0$）为
\[
  \mu([\ell_p,w];\lambda^0)=\mu(\mathcal V_w,\lambda_w).
\]
\end{definition}

由于我们不会改变参考截面 $\lambda^0$，以下把
$\mu([\ell_p,w];\lambda^0)$ 简记为 $\mu([\ell_p,w])$。关于弗洛尔链模的分次如何依赖于 $\lambda^0$ 的选取，见注 2.4.45。
